\documentclass[10pt,a4paper]{amsart}
\usepackage[margin=19mm]{geometry}
\usepackage{amsmath,amssymb,mathtools}
\usepackage[scr=rsfs,cal=euler]{mathalfa}
\usepackage{xfrac}
\usepackage[unicode,hidelinks,bookmarks=false]{hyperref}
\newcommand{\ie}{i.e.\ }
\newcommand{\cf}{cf.\ }
\newcommand{\resp}{respectively\ }
\newcommand{\Subsection}{\subsection}
\providecommand{\nobreakdash}{\nobreak\hbox{-}}
\setlength{\emergencystretch}{2em}
\multlinegap0pt
\usepackage{amssymb}
\usepackage{mathtools}
\usepackage{bm}
\usepackage[shortlabels]{enumitem}
\usepackage{stackengine}
\newtheorem{theorem}{Theorem}
\newtheorem{lemma}{Lemma}
\title{Liouville Theorems Above the Critical $9/2$ Threshold for Stationary Navier-Stokes Equations}
\author{Gaston Vergara-Hermosilla}
\date{Source: arXiv:2604.06527v2 (CC BY 4.0)}
\begin{document}
\maketitle

\noindent\emph{Source and licence.} Liouville Theorems Above the Critical $9/2$ Threshold for Stationary Navier-Stokes Equations. Version arXiv:2604.06527v2 (CC BY 4.0), distributed by arXiv under the Creative Commons Attribution 4.0 International licence, http://creativecommons.org/licenses/by/4.0/, as recorded in the arXiv licence metadata for this exact version and shown on the abstract page https://arxiv.org/abs/2604.06527v2. Full licence notice: \texttt{licenses/arxiv-2604.06527-CC-BY-4.0.txt}.\par\smallskip

\noindent\textit{Editorial scope.} Counted unit: the article's lemma lem:PCsup\_asymptotics (source0.tex lines 546--558). The article's complete original proof of it is delivered with it (source0.tex lines 559--619, one \texttt{proof} environment of 61 lines). No mathematics is written, repaired or re-derived: the statement and the proof are the article's own lines, in the article's own notation, and no retained line is substituted or re-flowed.  Retained uncounted as the source-local context the proof needs: lines 211--214; lines 288--302.  Nothing was omitted from any retained span, so the package carries no omission record.  No retained line contains a citation, so the wrapper carries no bibliography and no bibliography entry.  No retained line contains a figure environment, an image-inclusion command or drawing code, so no image is delivered and the package declares no asset dependency.  Editorial formatting only, in addition: an AMS \texttt{amsart} preamble that restates the article's own declarations and macros used by the retained lines, namely the environments \texttt{theorem}, \texttt{lemma} on separate counters, together with the standard packages those lines need.  The retained environments print in this document as Theorem 1, Lemma 1 in order of appearance, whereas the article prints the same items with the numbering of its own full document.  The complete native source of the article as distributed in the arXiv e-print for this exact version is bundled unchanged as \texttt{sources/198022/source0.tex}.
\begin{equation}\label{SNS}
-\Delta u + u \cdot \nabla u + \nabla P = 0, \quad 
\nabla \cdot u = 0,
\end{equation}

\begin{theorem}\label{thm:over92}
Let $R_0 > 1$ be fixed, and let $u \in \dot{H}^1(\mathbb{R}^3)$ be a solution of \eqref{SNS}. Let $p:\mathbb{R}^3 \to \mathbb{R}$ be a variable exponent such that:
\begin{enumerate}
\item $p(x) = 6$ for all $x \in B(0,R_0)$,
\item $p$ is 
Lipschitz,  
 %continuous,
 radially decreasing, and satisfies $p(x) \ge \frac{9}{2}$ for all $x \in \mathbb{R}^3$,
\item there exists a constant $C \in [0, R_0^2)$ such that
$
|p(x) - \tfrac{9}{2}| \le \frac{C}{|x|}, \  \text{for all } |x| \ge R_0^2.
$
\end{enumerate}
If, in addition, $u \in L^{p(\cdot)}(\mathbb{R}^3)$, then $u \equiv 0$.
\end{theorem}

\begin{lemma}\label{lem:PCsup_asymptotics}
Let $p(\cdot)$ be the variable exponent defined in Theorem \ref{thm:over92} and $R\geq 2R_0^2$. 
Then, 
there exists a sequence $(\varepsilon_R)$ and $C>0$ such that
\[
p_{\mathcal{C}_R}^+ = \frac{9}{2} + \varepsilon_R, 
\quad 0<\varepsilon_R \le \frac{2C}{R}<1.
\]
In particular,
$\label{O1R}
p_{\mathcal{C}_R}^+ =  {9}/ {2} + O\!\left( {1} / {R}\right).
$
\end{lemma}

\begin{proof}
Let $x\in \mathcal{C}_R$. 
Since $p(\cdot)$ is continuous and radially decreasing,  
there exists a decreasing and continuous function
$\tilde{p} : [0, \infty) \to [9/2, 6]$, 
such that $ p(x) = \tilde{p}(|x|) $ for all $ x \in \mathbb{R}^3 $.
Then, for each $ x \in \mathcal{C}_R $ we have $|x| > R/2$, hence, 
\[
p(x) = \tilde{p}(|x|) \leq \tilde{p}(R/2) %= p(R/2).
,
\]
and we conclude 
\[
\text{ess sup}_{x \in \mathcal{C}_R} p(x) \leq  \tilde{p}(R/2) 
.
\] 
Let $ \epsilon > 0 $. By continuity of $ \tilde{p} $ at $ R/2 $, there exists $ \delta > 0 $ such that, 
$
r \in (R/2, R/2 + \delta) $ implies $|\tilde{p}(r) - \tilde{p}(R/2)| < \epsilon
,
$
in particular,  if
$
r \in (R/2, R/2 + \delta) $, then  $\tilde{p}(r) >\tilde{p}(R/2) - \varepsilon.$
To continue, we define the set
\[
A_\delta := \{x \in \mathbb{R}^3 : R/2 < |x| < R/2 + \delta\}.
\]
At this point we must stress the fact that, $ A_\delta \subset \mathcal{C}_R $ and, for all $ x \in A_\delta $, and we can write
\[
p(x) > \tilde{p}(R/2) - \epsilon.
\]
Moreover, $ A_\delta $ has positive Lebesgue measure, since
\[
| A_\delta | 
=
 \frac{4\pi}{3} ((R/2 + \delta)^3 - (R/2)^3) > 0.
\]
As a consequence of this, we get that the set $
\{x \in \mathcal{C}_R : p(x) > \tilde{p}(R/2) - \varepsilon\}$ has positive measure. Then, considering the definition of essential supremum, we can write $\operatorname{ess\,sup}_{x \in \mathcal{C}_R} p(x) \geq \tilde{p}(R/2) - \epsilon$. Thus, since  $ \varepsilon > 0 $ is arbitrary, we conclude
\[
\operatorname{ess\,sup}_{x \in \mathcal{C}_R} p(x) \geq \tilde{p}(R/2),
\]
and then, 
\[
p_{\mathcal{C}_R}^+ = \text{ess sup}_{x \in \mathcal{C}_R} p(x) =\tilde{p}(R/2).
\]
Now, to conclude  the asymptotic estimate, note that  
since $ R \geq 2R_0^2 $, we have $ R/2 \geq R_0^2 $, and by hypothesis, we know that there exist $C<R^2_0\leq R/2$ such that 
\[
\left| \tilde p(R/2) - \frac{9}{2} \right| \leq \frac{C}{R/2} = \frac{2C}{R}<1.
\]
Thus, defining %, for instance,  
  $\epsilon_R := \tilde p(R/2) - \frac{9}{2}$ 
%we obtain
%\[
%p_{\mathcal{C}_R}^+ = \frac{9}{2} + \varepsilon_R, \quad |\varepsilon_R| \leq \frac{2C}{R}<1,
%\]
%and
 we deduce the desired results. 
\end{proof}

\end{document}
